%% kapa: exam=1 round=28 year=2017 period=1 form=A subject=civil_law subject_ko=민법 seq=35 number=35
%% source: Q-Net 감정평가사 기출문제 (cst00309) artlSeq=5199603 file=extracted/1008/5199603/2017년 제28회 감정평가사 1차시험 문제지/1교시 A형.pdf
%% conversion: kordoc 4.24.51 -> tools/gen1.py (auto); review=proofread-sonnet; flags=none
\begin{question}{35}
전세권에 관한 설명으로 옳은 것은? (다툼이 있으면 판례에 따름)
\begin{choices}
  \item 전세권자의 책임 없는 사유로 전세권의 목적물 전부가 멸실된 때에도 전세권자는 손해배상책임이 있다.
  \item 건물에 대한 전세권이 법정갱신되는 경우, 그 존속기간은 2년으로 본다.
  \item 전세권의 존속기간이 만료되면 전세권의 용익물권적 권능은 전세권설정등기의 말소 없이도 당연히 소멸한다.
  \item 전세권설정자는 특약이 없는 한 목적물의 현상을 유지하고 그 통상의 관리에 속한 수선을 해야 한다.
  \item 전세권을 목적으로 저당권을 설정한 자는 저당권자의 동의 없이 전세권설정자와 합의하여 전세권을 소멸시킬 수 있다.
\end{choices}
\end{question}
\answer{3}
